%==========================================================================
\section{Constraints on the individual sources}
\label{sec:results}
%==========================================================================

For each source of Table~\ref{tab:catalog} we search the $(a,M)$ plane for valid configurations, i.e.\ for a precession radius $\bar r$ that satisfies Eqs.~\eqref{eq2:Pcond} and \eqref{eq2:Ncond}, using the fiducial parameters $h_\textup{ISCO}=0.1$ and $\alpha=0.01$.
%The dependence on $h_\textup{ISCO}$ and $\alpha$ is discussed in Appendix~\ref{app:params}, and the regularities that emerge from the individual sources are generalised to the whole $(M,\nu_0)$ plane in Sect.~\ref{sec:general}.

The precession frequency of Eq.~\eqref{eq2:gdef} is defined in the rest frame of the source, whereas Table~\ref{tab:catalog} lists the observed frequencies.
Cosmological time dilation stretches every period by $1+z$, so we compare the model with
\begin{equation}
    \label{eq3:nurest}
    \nu_0=\nu_\textup{obs}\,(1+z),
    \qquad
    P=m\,\nu_0=m\,\nu_\textup{obs}\,(1+z).
\end{equation}
The correction increases $P$ by 0.049~dex for J1231 ($z=0.119$) and by less than 0.02~dex for all the other sources, much less than the uncertainty on the masses; it changes the bounds on the spin by at most 0.04.

As a first step we impose only the frequency condition, Eq.~\eqref{eq2:Pcond}, without the timescale conditions.
Figure~\ref{fig:solid_aR} shows, for each source, the outer radii at which the flow precesses at the observed frequency for some mass within the published range.
Two features are already apparent.
First, every band is bounded from below by the ISCO: at a given spin, the fastest precession is that of a ring at the ISCO, so for $|a|$ below a minimum value no radius reproduces $\nu_0$.
For the sources with $\nu_0\gtrsim10^{-4}$~Hz and $M\gtrsim10^6\,M_\odot$ this rules out retrograde spins and requires $r_\textup{out}\lesssim15\,R_g$.
Second, the sources with lower $P$ (J1347, J1257, Mrk~142, J1231 and NGC~4945) admit both spin signs and larger radii, up to several $10^2\,R_g$ for NGC~4945.
In the following we add the conditions of Eq.~\eqref{eq2:Ncond}, which select the part of each band where the flow precesses rigidly for at least one period.
\begin{figure}
    \centering
    \includegraphics[width=\linewidth]{3_sources/pif_solid_aR.pdf}
    \caption{
    Outer radii $r_\textup{out}$ at which the flow precesses at the rest-frame QPO frequency, $\nu_\textup{solid}(a,r_\textup{out})=\nu_\textup{obs}(1+z)$, as a function of spin, for each source of Table~\ref{tab:catalog} (split between the two panels in the order of the table).
    The shaded band of each source is bounded by the curves for the lowest and highest mass of its published range.
    The dashed line is $r_\textup{ISCO}(a)$.
    }
    \label{fig:solid_aR}
\end{figure}


%--------------------------------------------------------------------------
\subsection{Valid configurations}
\label{sec3:grids}

Figure~\ref{fig:grid} shows the valid configurations of each source in the $(a,M)$ plane, and Table~\ref{tab:constraints} summarises them within the published mass ranges.
Since validity depends on $M$ and $\nu_0$ only through $P$, the valid regions of all the sources have the same shape, shifted in $\log M$ by $-\log\nu_0$; the sources differ in where their mass range falls within this shape.

The shape has two edges.
The upper edge, at high mass, is a geometric limit.
By Eq.~\eqref{eq2:signLT}, $g_\textup{LT}$ has the same sign at all radii, so that $g(a,r_\textup{out})$ in Eq.~\eqref{eq2:gdef} is the $w$-weighted mean of $|g_\textup{LT}|$ over $[r_\textup{ISCO},r_\textup{out}]$.
Since $|g_\textup{LT}|$ decreases with $r$, this mean is lower than the value at the inner edge, and tends to it for $r_\textup{out}\to r_\textup{ISCO}$:
\begin{equation}
    \label{eq3:gmax}
    g(a,r_\textup{out})<g_\textup{max}(a)\doteq\left|g_\textup{LT}\left(r_\textup{ISCO}(a),a\right)\right| .
\end{equation}
The fastest flow is thus the narrowest one, a ring at the ISCO, which precesses at the test-particle frequency (Sect.~\ref{sec2:pif}).
Equation~\eqref{eq2:Pcond} therefore has a solution only if $P<g_\textup{max}(a)$, that is, at a given spin, only for
\begin{equation}
    \label{eq3:Mmax_a}
    M<g_\textup{max}(a)\,\frac{M_\odot}{\nu_0},
\end{equation}
whatever $h_\textup{ISCO}$ and $\alpha$.
For $a<1$ this is also the edge of the valid region, because for $\bar r\to r_\textup{ISCO}$ the wave-propagation time vanishes and the alignment time diverges, so that Eq.~\eqref{eq2:Ncond} holds.

The function $g_\textup{max}$ vanishes for $a=0$ and increases with $|a|$ for both signs of the spin, but much faster for prograde spins, whose ISCO is closer to the hole.
For retrograde spins it reaches $110$~Hz at $a=-1$, a value that prograde spins already attain at $a=0.25$.
For prograde spins it is $1.05\times10^4$~Hz at $a=0.998$ and tends to $c/(4\pi R_{g,\odot})=1.62\times10^4$~Hz for $a\to1$, where $r_\textup{ISCO}\to1$ and $\Omega_z(r_\textup{ISCO})\to0$, so that $g_\textup{LT}\to g_\varphi=c/[2\pi R_{g,\odot}(1+a)]$ (Eqs.~\ref{eq2:epiz} and \ref{eq2:scale_nu}).
Each source has therefore a maximal mass
\begin{equation}
    \label{eq3:Mmax}
    M_\textup{max}=1.62\times10^{4}\,M_\odot\left(\frac{\nu_0}{\textup{Hz}}\right)^{-1},
\end{equation}
and retrograde configurations exist only for
\begin{equation}
    \label{eq3:Mmaxretro}
    M<110\,M_\odot\left(\frac{\nu_0}{\textup{Hz}}\right)^{-1},
\end{equation}
i.e.\ for $P<110$~Hz\,$M_\odot$.
The limit of Eq.~\eqref{eq3:Mmax} is approached for $a\to1$ and $\bar r\to r_\textup{ISCO}$ but is not reached, because no configuration with $a=1$ is valid (Sect.~\ref{sec2:numerics}).
If the spin is limited to $a\leq0.998$ \citep{Thorne1974}, the coefficient of Eq.~\eqref{eq3:Mmax} decreases by a factor of 1.53, to $1.05\times10^4\,M_\odot$.

At lower masses the precession radius moves outwards.
The lower edge of the valid region is where it moves so far that the flow aligns within one period, or, for $a\gtrsim0.9$, where bending waves no longer cross it within one period.
Because the lower edge depends on $h_\textup{ISCO}$ and $\alpha$ and the upper one does not, the bounds on $|a|$ from below and on $M$ from above in Table~\ref{tab:constraints} are independent of the disc parameters, whereas the bounds on $|a|$ from above are not (Appendix~\ref{app:params}).

According to Eq.~\eqref{eq3:Mmaxretro}, the sources fall into two groups: those with $P>110$~Hz\,$M_\odot$ over their whole mass range, for which retrograde flows are excluded (Sect.~\ref{sec3:highP}), and those whose mass range extends below this value (Sect.~\ref{sec3:lowP}).

\begin{figure*}
    \centering
    \includegraphics[width=0.9\hsize]{3_sources/pif_timescales_parameter_grids.pdf}
    \caption{
    Valid configurations in the $(a,M)$ plane for each source of Table~\ref{tab:catalog}, for the fiducial parameters and the rest-frame frequency.
    Green: valid within the published mass range; grey: valid outside it.
    }
    \label{fig:grid}
\end{figure*}

\input{3_sources/pif_source_constraints.tex}
% generated by pif_constraints.py; edit there, not here
% NOTE: the table must be regenerated with P_MAX evaluated at a = 1 (see Eq. eq3:Mmax)

%--------------------------------------------------------------------------
\subsection{Fast prograde spins and compact flows}
\label{sec3:highP}

Ten of the fifteen sources (RE~J1034+396, MS~2254.9-3712, Mrk~766, MCG-06-30-15, ESO~113-G010, 1H~0707-495, NGC~1365, NGC~5506, NGC~4151 and 1ES~1927+654) have $P>110$~Hz\,$M_\odot$ over their whole mass range.
Retrograde flows are excluded for all of them, and a prograde spin above a minimum value is required, between $0.29$ (RE~J1034+396, at its lowest mass) and $0.84$ (1ES~1927+654 at 2.4~mHz).
The precessing flow is compact: $\bar r\lesssim9\,R_g$ (up to $11\,R_g$ at the lowest mass of RE~J1034+396), and in most cases $\bar r\simeq2$--$6\,R_g$, at most a few times $r_\textup{ISCO}$.
The sources with the narrowest mass ranges are the most constraining: the valid spins lie in $0.65\lesssim a\lesssim0.96$ for MS~2254.9-3712 and in $0.47\lesssim a\lesssim0.95$ for ESO~113-G010.

Figure~\ref{fig:mrk766} illustrates this group with Mrk~766, whose published masses range from $1.3\times10^6\,M_\odot$ \citep{Giacche2014} to the reverberation-mapping value of $6.6\times10^6\,M_\odot$ \citep{BentzKatz2015}, i.e.\ $P\simeq200$--$1200$~Hz\,$M_\odot$.
At the geometric mean of the range ($M=3.1\times10^6\,M_\odot$, left panel), no flow precesses fast enough for $a\leq0.5$: even a ring at the ISCO is slower than the QPO.
Of the spins shown, only $a=0.9$ is valid, with $\bar r=6.0\,R_g$.
At fixed $a=0.9$ (right panel) the precession radius moves inwards as the mass increases, from $6.6\,R_g$ at $2.6\times10^6\,M_\odot$ to $3.7\,R_g$ at $7.4\times10^6\,M_\odot$, and all these configurations are valid.
For $M\leq1.8\times10^6\,M_\odot$ it lies instead at $\bar r\geq7.9\,R_g$, where both conditions fail.
The required spin therefore increases with the mass: $0.39\leq a\leq0.85$ at the lowest mass and $0.76\leq a\leq0.94$ at the reverberation-mapping value.
For the highest masses the upper bound on the spin is set by wave propagation, because near the ISCO of a rapidly spinning hole $\Omega_z\ll\Omega_\varphi$ and the sound speed $c_s=H\Omega_z$ is small (Sect.~\ref{sec2:timescales}).

\begin{figure*}
    \centering
    \includegraphics[width=0.9\hsize]{3_sources/pif_timescales_Mrk_766.pdf}
    \caption{
    Wave-propagation time (dotted) and alignment time (dashed) as functions of $r_\textup{out}$ for Mrk~766 ($\nu_0=1.57\times10^{-4}$~Hz), for the fiducial parameters.
    Left: spins between $-1$ and $1$ at the geometric mean of the mass range, $M=3.1\times10^6\,M_\odot$.
    Right: masses spanning the range of Table~\ref{tab:catalog} at $a=0.9$.
    The horizontal line is the period $1/\nu_0$.
    Circles mark the precession radii $\nu_\textup{solid}(\bar r)=\nu_0$: filled if $t_\textup{wave}(\bar r)<1/\nu_0<t_\textup{align}(\bar r)$, open otherwise.
    Curves without a circle have no precession radius.
    For $a=1$, $t_\textup{wave}$ diverges and is not shown.
    }
    \label{fig:mrk766}
\end{figure*}

\paragraph{Upper limits on the mass.}
For two entries, NGC~4151 and the higher-frequency QPO of 1H~0707-495, $M_\textup{max}$ falls below the upper end of the mass range (daggers in Table~\ref{tab:constraints}), so that part of the published estimates is incompatible with the model for any spin and any $(h_\textup{ISCO},\alpha)$.
The most significant case is NGC~4151, the source with the highest $P$, for which $M_\textup{max}=2.7\times10^7\,M_\odot$.
The upper fifth of its mass range in logarithm, which includes the reverberation-mapping estimate \citep{Bentz2006}, is excluded, and the QPO requires both a mass in the lower part of the published range and $a\gtrsim0.83$.
For instance, at $a=0.95$ precession radii exist only for $M\lesssim6.5\times10^6\,M_\odot$, with $\bar r\simeq2.4$--$3.5\,R_g$.
For 1H~0707-495 (QPO a), $M_\textup{max}=5.9\times10^7\,M_\odot$ excludes the central value of the highest estimate in Table~\ref{tab:catalog} \citep{Titarchuk2025}, although not the lower end of its uncertainty.

These limits require a spin arbitrarily close to one.
For $a\leq0.998$ they decrease by a factor of 1.53: the limit for NGC~4151 becomes $1.8\times10^7\,M_\odot$, which excludes the upper third of its mass range in logarithm, and the limit for 1H~0707-495 (QPO a) becomes $3.9\times10^7\,M_\odot$, below the whole uncertainty interval of the highest estimate.
Four more entries are then affected.
For RE~J1034+396 and MCG-06-30-15 the limit, $\simeq3.8\times10^7\,M_\odot$, excludes the highest estimates in Table~\ref{tab:catalog} \citep{Czerny2016,Das2025}.
For NGC~1365, $4.8\times10^7\,M_\odot$ excludes the estimate from the bulge luminosity \citep{Risaliti2009}, and for 1ES~1927+654 at 2.4~mHz, $4.2\times10^6\,M_\odot$ excludes the highest of the host-galaxy estimates \citep{Li2022}.
For the second QPO of 1H~0707-495 the limit, $8.4\times10^7\,M_\odot$, is close to the top of the range.
In all these sources the lower estimates remain allowed, with $a\gtrsim0.3$--$0.85$.
% check that the excluded estimates are the ones attributed to these references in Table 1
% Optional figure: output_catalog/pif_timescales_NGC_4151.pdf (right panel at a = 0.95)
% shows the mass cut directly.

%--------------------------------------------------------------------------
\subsection{Both spin signs and larger flows}
\label{sec3:lowP}

For J1347, J1257, Mrk~142, J1231 and NGC~4945 the mass range extends below $P=110$~Hz\,$M_\odot$, so that retrograde flows are possible.
Every mass of their ranges is allowed for some prograde spin.
Retrograde spins are allowed over the whole range only for J1231 and NGC~4945 ($P\leq8$~Hz\,$M_\odot$).
For the other three sources, whose ranges reach $P\simeq200$--$340$~Hz\,$M_\odot$, they are restricted by Eq.~\eqref{eq3:Mmaxretro} to $M<3.3\times10^6\,M_\odot$ (J1257), $M<1.5\times10^6\,M_\odot$ (Mrk~142), and $M<9.0\times10^6\,M_\odot$ (J1347, higher frequency); at higher masses only prograde spins are allowed, as in Sect.~\ref{sec3:highP}.
The spin is less constrained than in the first group: for J1347, J1257 and Mrk~142, prograde spins between $\simeq0.1$--$0.2$ and $0.8$--$0.9$ and retrograde spins with $|a|\gtrsim0.1$--$0.4$ are valid for some mass of the range.
The precessing flow is larger, $\bar r\simeq4$--$20\,R_g$.

Figure~\ref{fig:j1257} shows the timescales of J1257 at the geometric mean of its mass range, $M=2.4\times10^6\,M_\odot$ ($P=82$~Hz\,$M_\odot$), close to the virial estimates.
Retrograde spins with $|a|\geq0.5$ and prograde spins $0.25\leq a\leq0.5$ are valid, with $\bar r\simeq6$--$11\,R_g$.
For $a=-0.25$ and $|a|\leq0.01$ even a ring at the ISCO precesses too slowly, while for $a\geq0.9$ the precession radius moves out to $15$--$21\,R_g$, where both conditions fail.
At fixed $a=0.5$ (right panel) the whole mass range is valid, and the precession radius decreases from $15\,R_g$ at $10^6\,M_\odot$ to $5.8\,R_g$ at $5.9\times10^6\,M_\odot$.
Over the mass range, J1257 requires $0.11\leq a\leq0.84$ or $a\leq-0.14$.

\begin{figure*}
    \centering
    \includegraphics[width=0.9\hsize]{3_sources/pif_timescales_J1257.pdf}
    \caption{
    As Fig.~\ref{fig:mrk766}, for J1257 ($\nu_0=3.37\times10^{-5}$~Hz).
    Left: spins between $-1$ and $1$ at the geometric mean of the mass range, $M=2.4\times10^6\,M_\odot$.
    Right: masses spanning the range of Table~\ref{tab:catalog} at $a=0.5$.
    }
    \label{fig:j1257}
\end{figure*}

Two sources require a slowly spinning hole.
J1231, the least massive of the sample ($3$--$10\times10^4\,M_\odot$, $P=2.5$--$8$~Hz\,$M_\odot$), is valid only for $0<a\leq0.38$ or $a<0$, with the largest precession radii of the sample, up to $46\,R_g$ (Fig.~\ref{fig:2xmm}).
At these radii $t_\textup{wave}$ and $t_\textup{align}$ are both close to $1/\nu_0$, so that the range of valid radii is narrow.
For $a=0.5$ no mass of the range is valid, because the flow aligns within one period at all the precession radii ($31$--$58\,R_g$); for $a=0.9$ both conditions fail.
NGC~4945, with the lowest frequency of the sample ($P=0.2$--$2$~Hz\,$M_\odot$), requires $0<a\leq0.21$ or $-0.39\leq a<0$, and $\bar r\simeq6$--$74\,R_g$, partly outside the bending-wave regime (Sect.~\ref{sec2:timescales}).
At the mass given by the rotation of the masers, $1.4\times10^6\,M_\odot$ \citep{Greenhill1997}, the valid spins are $|a|\lesssim0.1$.

\begin{figure*}
    \centering
    \includegraphics[width=0.9\hsize]{3_sources/pif_timescales_J1231.pdf}
    \caption{
    As Fig.~\ref{fig:mrk766}, for J1231 ($\nu_0=8.18\times10^{-5}$~Hz).
    Left: spins between $-1$ and $1$ at the geometric mean of the mass range, $M=5.5\times10^4\,M_\odot$.
    Right: masses spanning the range of Table~\ref{tab:catalog} at $a=0.5$.
    }
    \label{fig:2xmm}
\end{figure*}

%--------------------------------------------------------------------------
\subsection{Sources with two QPOs}
\label{sec3:pairs}

Four sources show two frequencies (Table~\ref{tab:catalog}), which must be reproduced by the same $M$ and $a$.

\paragraph{1ES~1927+654.}
The two frequencies were observed at different epochs, as the QPO drifted from 0.9 to 2.4~mHz \citep{Masterson2025}.
In the model this is a single flow around a fixed BH whose outer radius shrinks.
Both frequencies are valid together only for
\begin{equation}
    \label{eq3:1es}
    0.85\lesssim a\lesssim0.96,\qquad M\lesssim1.6\times10^6\,M_\odot ,
\end{equation}
which excludes the upper part of the published mass range (Table~\ref{tab:catalog}).
Within this region the precession radius decreases from $\bar r\simeq3$--$5\,R_g$ at 0.9~mHz to $\bar r\simeq2$--$3\,R_g$ at 2.4~mHz: the drift corresponds to a contraction of the precessing region by a factor of $\simeq1.8$, after which it remained at $\bar r\simeq2$--$3\,R_g$ as the frequency stabilised \citep{Masterson2026}.

\paragraph{Simultaneous pairs.}
In ESO~113-G010, 1H~0707-495 and J1347 the two QPOs were detected in the same data, with frequency ratios of 1.8, 2.2 and 3.0.
A joint solution exists for all three: $0.61\lesssim a\lesssim0.93$ for ESO~113-G010 over its whole mass range, $0.62\lesssim a\lesssim0.97$ and $M\lesssim1.7\times10^7\,M_\odot$ for 1H~0707-495, and $-1\leq a\lesssim0.78$ for J1347.
A joint solution, however, assigns two different outer radii to the same flow at the same time, whereas a single rigidly precessing flow has a single precession frequency.
If the model applies, the two peaks of each pair therefore cannot be two independent precession frequencies.
They can instead be interpreted as harmonics of the same precession: a periodic modulation that is not sinusoidal has power not only at its fundamental frequency, the precession frequency, but also at integer multiples of it.
This is commonly observed for type-C QPOs in XRBs, whose power spectra show a second harmonic, and often a subharmonic, besides the fundamental \citep{ingram}.
In this interpretation only the fundamental is the precession frequency, and the constraints on the source are those of the single entry of Table~\ref{tab:constraints} that is identified with it.

%--------------------------------------------------------------------------
\subsection{Summary}
\label{sec3:summary}

For the fiducial parameters, every source of the catalogue admits valid configurations within its published mass range, but only in a restricted part of the $(a,M)$ plane:
\begin{itemize}
    \item the ten sources with $P>110$~Hz\,$M_\odot$ over their whole mass range require a prograde spin, $a\gtrsim0.3$--$0.85$, and a compact precessing flow, $\bar r\lesssim9\,R_g$;
    \item the mass is bounded from above by Eq.~\eqref{eq3:Mmax}, independently of the disc parameters, which excludes part of the published estimates for NGC~4151 and 1H~0707-495 whatever the spin, and also for RE~J1034+396, MCG-06-30-15, NGC~1365 and 1ES~1927+654 if $a\leq0.998$;
    \item the five sources whose mass range extends below $P=110$~Hz\,$M_\odot$ allow larger flows, $\bar r\simeq4$--$75\,R_g$, and retrograde spins for $M<110\,M_\odot\,(\nu_0/\textup{Hz})^{-1}$; the two with the lowest $P$ (J1231, NGC~4945) require $a\lesssim0.2$--$0.4$ if the flow is prograde.
\end{itemize}
The common shape of the valid regions in Fig.~\ref{fig:grid} reflects the fact that validity depends on $M$ and $\nu_0$ only through $P$.
In Sect.~\ref{sec:general} we use this property to extend the analysis from the observed sources to the whole $(M,\nu_0)$ plane.